\documentclass[../../../include/open-logic-section]{subfiles}

\begin{document}

\olfileid[hi]{sth}{cardinals}{hp}
\olsection{परिशिष्ट: ह्यूम का सिद्धांत}

\olref[cp]{sec} में हमने कैंटर के सिद्धांत का वर्णन किया था। वह था:
\begin{align*}
	\card{A} = \card{B} & \text{ तभी जब } A \approx B.
\intertext{यह उस सिद्धांत से बहुत मिलता-जुलता है जिसे अब
\emph{ह्यूम का सिद्धांत} कहते हैं और जो कहता है:}
	\fregenum{x} {F(x)} = \fregenum{x}{G(x)} & \text{ तभी जब } F \sim G
\end{align*}
जहाँ `$F\sim G$' इसका संक्षिप्त रूप है कि $F$ उतने ही हैं जितने $G$,
अर्थात् $F$ को $G$ के साथ एक द्विआक्षेपी संगति में रखा जा सकता है, अर्थात्:
\begin{align*}
	\exists R(&\forall v\forall y(Rvy \lif (Fv \land Gy)) \land {}\\
		&\forall v(Fv \lif \lexists![y][Rvy]) \land {}\\
		&\forall y(Gy \lif \lexists![v][Rvy]))
\end{align*}
किंतु ह्यूम के सिद्धांत और कैंटर के सिद्धांत के बीच प्रकार का अंतर है।
कैंटर के सिद्धांत के कथन में चर ``$A$'' और ``$B$'' प्रथम-क्रम पद हैं जो
\emph{समुच्चयों} के लिए खड़े होते हैं। ह्यूम के सिद्धांत के कथन में ``$F$'',
``$G$'' और ``$R$'' प्रथम-क्रम पद \emph{नहीं} हैं; बल्कि वे \emph{विधेय
स्थान} में हैं। (शायद वे \emph{गुणों} के लिए खड़े होते हैं।) अतः ह्यूम के
सिद्धांत को हम इस प्रकार कह सकते हैं: $F$ की संख्या $G$ की संख्या के बराबर
है तभी जब $F$ और $G$ के बीच द्विआक्षेपी संगति हो। इसे \emph{ह्यूम का
सिद्धांत} कहते हैं, क्योंकि ह्यूम ने एक बार यह लिखा था:
\begin{quote}
  जब दो संख्याएँ इस प्रकार संबद्ध हों कि एक की प्रत्येक इकाई के अनुरूप
  दूसरी की सदैव एक इकाई हो, तो हम उन्हें बराबर कहते हैं।
  \citep[Pt.III Bk.1 \S1]{Hume1740}
\end{quote}
और \citet[\S63]{Frege1884} ने ह्यूम के इस अंश को उद्धृत करके, फिर (वस्तुतः)
(जिसे हमने कहा है) ह्यूम के सिद्धांत की रूपरेखा प्रस्तुत करके, इस सिद्धांत
को समकालीन गणितीय-तार्किक महत्त्व दिलाया।

ह्यूम के सिद्धांत और मूल नियम~V के बीच संरचनात्मक समानता पर ध्यान दीजिए।
हमने इसे \olref[story][blv]{sec} में इस प्रकार सूत्रबद्ध किया था:
\[
	\fregeext{x}{F(x)} = \fregeext{x}{G(x)} \text{iff } \lforall[x][(F(x) \liff G(x))].
\]
इस बिंदु पर कुछ टिप्पणी और तुलना उपयोगी हो सकती है।

ह्यूम के सिद्धांत या मूल नियम~V जैसे सिद्धांत को समझने के दो ढंग हैं:
\emph{विधेयात्मक} या \emph{अविधेयात्मक}
(\olref[story][predicative]{sec} याद करें)। मूल नियम~V के अविधेयात्मक पाठ
में प्रत्येक $F$ के लिए वस्तु $\fregeext{x}{F(x)}$ उसी परिमाणीकरण-क्षेत्र
में आती है जिसका उपयोग हमने स्वयं मूल नियम~V को सूत्रबद्ध करने में किया था।
इसी प्रकार ह्यूम के सिद्धांत के अविधेयात्मक पाठ में प्रत्येक $F$ के लिए
वस्तु $\fregenum{x}{F(x)}$ उसी परिमाणीकरण-क्षेत्र में आती है जिसका उपयोग
हमने ह्यूम के सिद्धांत को सूत्रबद्ध करने में किया था। इसके विपरीत
\emph{विधेयात्मक} समझ में वस्तुएँ $\fregeext{x}{F(x)}$ और
$\fregenum{x}{F(x)}$ किसी \emph{भिन्न} क्षेत्र की सत्ताएँ होंगी।

यदि हम मूल नियम~V को अविधेयात्मक रूप से पढ़ें, तो सहज समुच्चय-निर्माण के
माध्यम से वह असंगति की ओर ले जाता है (विवरण के लिए
\olref[story][blv]{sec} देखें)। सहज समुच्चय-निर्माण की तरह ही उसे
\emph{विधेयात्मक} रूप से पढ़कर संगत बनाया जा सकता है। किंतु तब वह संभवतः वह
सब नहीं करेगा जो हम उससे चाहते थे।

किंतु ह्यूम के सिद्धांत को अविधेयात्मक रूप से संगत ढंग से पढ़ा \emph{जा
सकता है}। और इस प्रकार पढ़े जाने पर वह काफ़ी शक्तिशाली है।

उदाहरण के लिए विधेय ``$x\neq x$'' पर विचार करें, जिसे स्पष्टतः कोई वस्तु
संतुष्ट नहीं करती। ह्यूम का सिद्धांत अब एक वस्तु $\#x(x\neq x)$ देता है।
हम इसे संख्या~$0$ मान सकते हैं। अब \emph{अविधेयात्मक} समझ में---किंतु
\emph{केवल} अविधेयात्मक समझ में---यह सत्ता $0$ हमारे मूल परिमाणीकरण-क्षेत्र
में आती है। इसलिए हम विधेय ``$x=0$'' के साथ ह्यूम के सिद्धांत को सार्थक रूप
से लागू करके एक वस्तु $\#x(x=0)$ प्राप्त कर सकते हैं। हम इसे संख्या~$1$
मान सकते हैं। इसके अतिरिक्त ह्यूम के सिद्धांत से $0\neq1$ निष्पन्न होता है,
क्योंकि अपने से असमान वस्तुओं से $0$ के समान वस्तुओं तक कोई द्विआक्षेप नहीं
हो सकता (पहली कोई नहीं, किंतु दूसरी एक है)। अब फिर अविधेयात्मक ढंग से काम
करने पर $1$ हमारे मूल परिमाणीकरण-क्षेत्र में आता है। अतः हम विधेय
``$(x=0\lor x=1)$'' के साथ ह्यूम के सिद्धांत को सार्थक रूप से लागू करके
वस्तु $\#x(x=0\lor x=1)$ प्राप्त कर सकते हैं। हम इसे संख्या~$2$ मान सकते
हैं, और दिखा सकते हैं कि $0\neq2$ और $1\neq2$, और इसी प्रकार आगे।

संक्षेप में, अविधेयात्मक रूप से लेने पर ह्यूम के सिद्धांत से निष्पन्न होता
है कि \emph{अनंततः अनेक वस्तुएँ हैं}। इसने \emph{नव-फ़्रेगेय
तर्कवादियों} को ह्यूम के सिद्धांत को अंकगणित का आधार मानने के लिए प्रेरित
किया है।

किंतु फ़्रेगे ने \emph{स्वयं} ह्यूम के सिद्धांत को अंकगणित का आधार नहीं
माना। इसके बदले फ़्रेगे ने एक स्पष्ट परिभाषा से ह्यूम का सिद्धांत सिद्ध
किया: $\fregenum{x}{F(x)}$ को संकल्पना $F\sim\Phi$ के विस्तार के रूप में
परिभाषित किया गया है। आधुनिक शब्दों में हम इसे
$\fregenum{x}{F(x)}=\Setabs{G}{F\sim G}$ के रूप में प्रस्तुत करने का प्रयास
कर सकते हैं; किंतु यह हमें सहज समुच्चय-निर्माण की समस्याओं में वापस खींच
लेगा।

\end{document}
